\documentclass[runningheads]{llncs}
\usepackage[T1]{fontenc}
\usepackage{amsmath,amssymb}
\usepackage{graphicx}
\usepackage{booktabs}
\usepackage{algorithm}
\usepackage{algpseudocode}
\usepackage{url}
\usepackage{tikz}
\usetikzlibrary{arrows.meta,positioning}

\begin{document}

\title{Worked Examples That Compute Themselves: Step-by-Step Explainers and Generated Interactive Labs for Lecture Decks}
\titlerunning{Computed Explainers and Generated Labs for Lecture Decks}

\author{Swayam Devendra Panchal \and Kawaljeet Singh Bharaj \and Jeevesh Bharat Bodhani \and Rishit Pankajkumar Lahoti \and Pratik Kanani}
\authorrunning{S. D. Panchal et al.}
\institute{Department of Artificial Intelligence and Data Science,\\
Dwarkadas J. Sanghvi College of Engineering, Mumbai, India\\
\email{panchalswayam35@gmail.com}}

\maketitle

\begin{abstract}
Lecture slides in technical courses often contain a solved problem---a Viterbi trellis, a CYK chart, a Bayes calculation---
shown as a finished table. Students see the answer but not the order of decisions that produced it. We present two
components of Learnova, a system that rebuilds existing lecture decks, which turn such slides into something a student can
work through. The first is a set of exact solvers that recognise a problem in slide text and produce a beginner-level
explanation, one step at a time, with every number computed symbolically rather than generated by a language model. The
second is a lab generator that attaches an interactive step-through page to each topic of a deck: hand-built labs for
common topics, and labs generated for other topics either by a language model, whose code is
checked before use, or offline from the slide's own structure. On a held-out set of 26 problems written for this study the
solvers detected all problems and gave the correct answer in every case after one parsing fix, with no false detections on
12 ordinary lecture sentences and a median time of 0.7\,ms. On three real NLP decks the system attached 6, 11 and 14 labs
without any API key.
\keywords{worked examples \and step-by-step solutions \and interactive simulations \and NLP education \and
presentation generation}
\end{abstract}

\section{Introduction}
Worked examples are one of the best-supported techniques in instructional research: novices learn more from studying a
solved problem step by step than from solving the same problem unaided~\cite{sweller1985,atkinson2000,renkl2014}. The
benefit depends on two things that a lecture slide usually lacks. Each step must be visible on its own, and the reason for
the step must be stated. A slide that shows a completed Viterbi trellis gives neither.

Our course decks contain many such slides. When we
rebuilt them with Learnova we wanted every solved problem to become a stepper that a student can move
through, and every hard topic to have a small laboratory where the student can change the input and watch the algorithm
run. Two constraints shaped the design. First, the numbers must be right. Language models make arithmetic errors and
confident mistakes~\cite{ji2023,kasneci2023}, and an explanation that is wrong in one step does more harm than no
explanation. Second, the system must work in a classroom without network access or API credit.

The contributions of this paper are: (i) a detector and a library of exact solvers that turn a problem in slide text into
an explained, step-by-step solution, including NLP algorithms; (ii) a lab bench and a lab selection policy that attach one
interactive page per topic; (iii) a lab generator with two paths---model-written code that is statically checked, and an
offline builder that uses the slide's structure; and (iv) an evaluation on a held-out problem set and on real decks.

\section{Related Work}
Step-by-step solutions are offered by computer algebra services such as Wolfram$|$Alpha~\cite{wolframalpha}, whose
explanations target the mathematics rather than a specific lecture. Intelligent tutoring systems provide step-level
feedback and approach the effectiveness of human tutoring~\cite{vanlehn2011}, but they require authoring for each domain.
Algorithm visualisation has a long history in computing education; meta-analyses report that the learner's engagement
with the visualisation, not the animation itself, predicts learning~\cite{hundhausen2002,naps2002}. PhET simulations show
how inquiry-oriented interactive pages can be used at scale~\cite{wieman2008}. The ICAP framework~\cite{chi2014} and the
broader evidence for active learning~\cite{freeman2014} support moving students from passive viewing to interactive
work, and retrieval practice~\cite{roediger2006} motivates the ``predict before you reveal'' prompts in our labs.

\section{Computed Step-by-Step Explainers}
\subsection{Representation}
An explainer is a tuple $E = (\kappa, p, \langle s_1,\dots,s_n\rangle, a, \rho)$: the problem type $\kappa$, the problem
statement $p$, a sequence of steps, the answer $a$ and a short recap $\rho$. Each step
$s_i = (\mathit{title}_i, \mathit{work}_i, \mathit{why}_i, \mathit{rule}_i, G_i)$ has a short action title, one or more
lines of working (with the part that changed marked for highlighting), a plain-language reason, the general rule behind
the step and, for table algorithms, the current state of the table $G_i$ with the cells to highlight. An explainer is
accepted only if $n \ge 3$ and every step has both a title and a reason.

\subsection{Detection}
Slide text rarely contains a clean problem statement. The detector therefore builds candidate strings in order: the whole
text (for grammars and systems of equations), neighbouring pairs of lines that both contain ``$=$'', each line and each
sentence. Each candidate is offered to the solvers in a fixed order and the first accepted explainer wins:
\begin{equation}
E^{*} = \text{first}\{\, S_j(c) \ne \varnothing : c \in \mathit{Cand}(x),\ S_j \in \mathcal{S} \,\}.
\end{equation}
For algebra, the equation is extracted from the sentence by removing leading words (``Solve'', ``Find $x$ if'', ``Solve for
$y$ given'') and trailing clauses, and is rejected if either side still contains a word of two or more letters. This
conservative rule is what keeps ordinary lecture sentences from being treated as problems.

\subsection{Solvers}
All arithmetic uses SymPy~\cite{meurer2017} or exact fractions. Table~\ref{tab:solvers} lists the solvers. We describe the
NLP solvers, which are the main reason for this work.

\begin{table}[t]
\centering
\caption{Solvers and what each step shows}
\label{tab:solvers}
\begin{tabular}{p{3.4cm}p{8.2cm}}
\toprule
Solver & Steps \\
\midrule
Linear equation & write, expand, collect unknowns, remove constant, divide, check by substitution \\
Quadratic & standard form, factor or formula, both roots, check \\
Pair of linear equations & elimination with the multiplier shown, back-substitution, check \\
Percent, interest, speed, mean, median, probability, Pythagoras & formula, substitution, result with units \\
Bayes' rule & name the quantities, numerator, denominator, posterior \\
Derivative, integral & rule applied term by term \\
CYK parsing & chart cell by cell, with the rules that fire \\
Bigram probability & counts, each conditional probability, product \\
Naive Bayes (WSD) & priors, smoothed likelihoods, scores, decision \\
Viterbi decoding & initialisation, recursion per time step, back-pointers, path \\
\bottomrule
\end{tabular}
\end{table}

\textit{Bigram model.} For a corpus with sentence markers, the maximum-likelihood estimate~\cite{jurafsky2025}
\begin{equation}
P(w_i \mid w_{i-1}) = \frac{C(w_{i-1} w_i)}{C(w_{i-1})}
\end{equation}
is computed for each pair in the query sentence, and the sentence probability is the product
$P(w_{1:n}) = \prod_{i} P(w_i \mid w_{i-1})$, shown as an exact fraction and as a decimal.

\textit{Naive Bayes word sense disambiguation.} For senses $s$ with training contexts, vocabulary size $V$, and $N_s$ the
number of word tokens labelled $s$:
\begin{equation}
\hat{s} = \arg\max_{s} \; P(s) \prod_{j} P(f_j \mid s), \qquad
P(f_j \mid s) = \frac{C(f_j, s) + 1}{N_s + V},
\end{equation}
where $P(s)$ is the fraction of training contexts with sense $s$. Words of the test context that never occur in training
are listed and skipped, which is the convention used in our course.

\textit{Viterbi decoding.} For an HMM with start probabilities $\pi_j$, transitions $a_{ij}$ and emissions $b_j(o)$, the
solver fills
\begin{equation}
v_1(j) = \pi_j\, b_j(o_1), \qquad
v_t(j) = \max_{i} \, v_{t-1}(i)\, a_{ij}\, b_j(o_t), \qquad
\psi_t(j) = \arg\max_{i}\, v_{t-1}(i)\, a_{ij},
\end{equation}
showing every candidate product for every cell, and then follows the back-pointers $\psi$ from
$\arg\max_j v_T(j)$.

\textit{CYK parsing.} For a grammar in Chomsky normal form and words $w_1\dots w_n$, cell $T[i,j]$ holds every
non-terminal that derives $w_i\dots w_j$:
\begin{equation}
T[i,i] = \{A : A \to w_i\}, \qquad
T[i,j] = \bigcup_{k=i}^{j-1}\{A : A \to B\,C,\ B \in T[i,k],\ C \in T[k{+}1,j]\},
\end{equation}
and the sentence is accepted when $S \in T[1,n]$. Each step of the explainer is one cell, with the split points tried.

\textit{Language-model fallback.} With an API key, a problem that no solver recognises but that looks like a problem (it
contains a number or a verb such as ``solve'' or ``compute'') may be explained by a language model. Such an explainer is
labelled on the page as not computed by the system, has at most eight steps and is accepted only with at least three.

\section{Interactive Labs}
\subsection{The lab bench}
All labs share one page structure: presets, a view of the current state, a ``what is happening'' panel with
\emph{Step}, \emph{Back} and \emph{Play}, and a ``try your own input'' area. A lab supplies only its inputs and a function
$\mathrm{build}(x) \to \langle f_1,\dots,f_m \rangle$ that returns frames; each frame $f_k$ is one small decision with its
explanation and the view to show. This makes a new lab a single function, and it is also the contract we give to a
language model when it writes a lab.

\subsection{Choosing labs for a deck}
Each built-in lab declares a list of topic phrases. A slide gets a built-in lab when its title or text matches one of them.
Otherwise a lab may be generated, subject to a policy that keeps decks usable: one lab per topic (continuation slides share
it), the same lab at most twice per deck, at most ten generated labs per deck, and no lab for agenda, summary or quiz slides
or for titles that look like OCR debris (fewer than 70\% letters, unbalanced brackets, or no word of four letters).

\subsection{Generated labs}
\textit{Model-written.} With a key, the model receives the topic and the slide text and must return JSON with a title,
input fields and JavaScript defining \texttt{PRESETS}, \texttt{fill}, \texttt{readInputs} and \texttt{build}. The code is
rejected if it contains \texttt{fetch}, \texttt{XMLHttpRequest}, dynamic \texttt{import}, \texttt{eval}, cookies or local
storage, or a closing script tag, and it must pass \texttt{node --check}. A lab never talks to the network.

\textit{Offline.} Without a key, or when the model's lab is rejected, the lab is built from the visual structure the
decision engine found on the slide: a process is walked step by step, a comparison is explored criterion by criterion, a
timeline event by event, and a list of concepts is revealed one at a time after a ``think first'' prompt. A slide with fewer
than three structured items gets no generated lab.

Labs are generated in parallel (four workers) before the deck is built and cached by a hash of their content, so building
the same deck again costs nothing.

\section{Evaluation}
\subsection{Explainer accuracy}
We wrote 26 problems that do not appear in our unit tests---new numbers and new wordings across all solver types---and
worked out each answer by hand. We also collected 12 ordinary sentences from our lecture notes (definitions, history,
descriptions of algorithms) that must not be treated as problems. During this exercise two of our hand answers turned out
to be wrong (a bigram product and a Viterbi path); the solvers were right in both cases.

\begin{table}[t]
\centering
\caption{Held-out explainer benchmark}
\label{tab:bench}
\begin{tabular}{lc}
\toprule
Measure & Result \\
\midrule
Problems detected (initial run) & 25 / 26 \\
Problems detected (after one parsing fix) & 26 / 26 \\
Correct problem type & 26 / 26 \\
Correct answer & 26 / 26 \\
Every step has a reason & 26 / 26 \\
Mean steps per explainer & 5.1 \\
Median solve time & 0.7 ms \\
False detections on lecture sentences & 0 / 12 \\
\bottomrule
\end{tabular}
\end{table}

The one initial miss was the phrasing ``Find $x$ if $4x + 9 = x + 21$''; the extractor did not remove the words before the
equation. We added the pattern and two regression tests. We report both numbers because the fix was made after seeing the
benchmark.

\subsection{Labs on real decks}
On three NLP lecture decks (36, 95 and 216 pages), built without an API key, the system attached labs to 6, 11 and 14 distinct
topics. The 216-page syntax deck received four built-in labs (CYK parsing, Viterbi, bigram model, HMM tagging) and ten
generated ones, which is the cap. With API keys, generating labs sequentially had made a deck take 163\,s; generating them in
parallel before the build brought it to 74\,s for the 36-page deck.

\section{Discussion}
Exact solvers cover a closed set of problem types, and this is a deliberate trade. Within that set the explanations are
always correct and are produced in under a millisecond; outside it the system says honestly that the explanation came from
a model, or shows nothing. The labs follow the same principle: a generated lab that cannot be checked is not shown.

Our evaluation measures correctness and coverage, not learning. The research on worked examples and on interactive
visualisation~\cite{renkl2014,hundhausen2002} suggests that the format can help, but whether these particular steppers and
labs improve examination performance needs a classroom study, which we plan for the next semester.

\section{Conclusion}
Turning a solved problem on a slide into an explained, computed stepper, and a hard topic into a small lab, makes lecture
decks something students can work through rather than only read. By computing every number and checking every generated
program, the system can offer these features without the errors that make AI-generated study material hard to trust.

\begin{credits}
\subsubsection{\ackname} We thank our project mentor and the Department of Artificial Intelligence and Data Science, DJSCE,
for guidance and for the lecture material used in this study.
\subsubsection{\discintname} The authors have no competing interests to declare that are relevant to the content of this
article.
\end{credits}

\bibliographystyle{splncs04}
\bibliography{../refs}

\end{document}
